\section{神经网络的生物学启发}

\subsection{生物神经元 vs 人工神经元}

\begin{figure}[H]
\centering
\includegraphics[width=0.95\textwidth]{slides-images/slide-040.jpg}
\caption{生物神经元结构：树突（dendrites）接收信号，细胞体（cell body）汇聚处理，轴突（axon）传递信号\protect\footnotemark}
\end{figure}
\footnotetext{Slides 第41页。}

人工神经网络与生物神经元有松散的类比关系：

\begin{itemize}
    \item \textbf{树突}（dendrites）$\leftrightarrow$ 输入连接：接收来自其他神经元的信号
    \item \textbf{细胞体}（cell body）$\leftrightarrow$ 激活函数：汇聚输入信号并进行非线性变换
    \item \textbf{轴突}（axon）$\leftrightarrow$ 输出连接：将处理后的信号传递给下一层神经元
\end{itemize}

\begin{figure}[H]
\centering
\includegraphics[width=0.95\textwidth]{slides-images/slide-041.jpg}
\caption{人工神经元模型：输入加权求和后经过激活函数产生输出\protect\footnotemark}
\end{figure}
\footnotetext{Slides 第42页。}

\begin{warningbox}{脑科学类比需谨慎}
人工神经网络与生物神经系统的相似性非常\textbf{松散}。生物神经元远比人工神经元复杂：它们有不同类型的突触、时间动态、化学信号、稀疏连接等。人工神经网络为了计算效率使用规则的层状结构，而大脑的连接模式复杂得多。不应过度解读这种类比。
\end{warningbox}

\subsection{本章小结}

人工神经网络受到生物神经系统的启发，但已经发展出独立的理论和实践体系。现代深度学习的成功更多依赖于数学和工程优化，而非神经科学。

%% ========== 计算图 ==========

